\section{A Quantitative Spatial Model of Land-Use Approval}
\label{sec:model}

The reduced-form evidence shows that, after land-use authority shifts toward district-based control, large multifamily production shifts away from homeowner-heavy Council districts.
This quantity response is not by itself a welfare result.
Lower production in homeowner-heavy districts could reflect inefficient local veto power, but it could also reflect local officials screening out projects with high neighborhood costs.
To separate these interpretations, I outline a within-city quantitative spatial model with an institution-dependent land-use approval stage.

The household and firm side of the model follows the standard quantitative spatial tradition: households sort across locations, firms operate in locations with different productivity, and rents, wages, commuting flows, and housing supply clear in equilibrium.
The additional object is the land-use approval stage.
Developers may find housing projects privately feasible, but the amount of housing that actually enters the spatial equilibrium depends on whether those projects clear the political approval process.

The political structure borrows from the urban-growth logic in \citet{duranton_urban_2023}.
Incumbents may favor restrictions when they weigh agglomeration or affordability benefits against local urban costs, but they need not internalize the welfare of excluded newcomers.
I adapt this logic to a within-city setting.
District-level decision makers may use valuable local information about neighborhood costs, but they may also overweight incumbent homeowners and ignore benefits to entrants and other districts.
Centralized decision makers may internalize broader benefits, but may mismeasure true local costs.

\subsection{Overview and timing}

The model has three blocks.
First, given approved housing, a standard spatial equilibrium determines household location choices, rents, wages, entry, and asset values.
Second, developers face physical construction costs and a political or regulatory approval wedge.
Third, the approval regime determines approved housing.
The decentralized regime maximizes a politically weighted incumbent objective and observes local costs; the centralized regime maximizes a broader citywide objective but observes local costs imperfectly.

The timing is as follows.
Locations have baseline fundamentals and predetermined homeowner exposure.
Developers face market returns and physical redevelopment costs.
The approval regime determines which privately feasible housing clears the political process.
Households and firms then sort across the resulting spatial equilibrium.
For notational simplicity I suppress time in the model exposition.
Empirically, approved housing should be interpreted as new approved or built housing over a period, corresponding to the flow outcome used in the reduced form.

\subsection{Environment}

The city consists of neighborhoods indexed by \(n\in\mathcal{N}\).
Each neighborhood has baseline housing stock \(\bar H_n\), approved new housing \(q_n\), total housing
\[
    H_n = \bar H_n + q_n,
\]
residential rents or floor-space prices \(Q_n\), and predetermined homeowner exposure \(T_n\).
The object \(q_n\) is approved and built housing, not proposed housing.
This distinction matters because the political approval process is the margin the model is designed to capture.

Locations are stamped with baseline physical and amenity fundamentals: existing housing, developable capacity, construction costs, residential amenities, commuting access, and true local costs of additional development.
Rents, wages, population, home values, and approved housing are equilibrium outcomes.
Homeowner exposure \(T_n\) is treated as a predetermined political exposure measure.

There are three household groups:
\[
    g\in\{O,R,N\},
\]
where \(O\) denotes incumbent homeowners, \(R\) denotes incumbent renters, and \(N\) denotes potential entrants.
Homeowners hold local housing assets and may gain or lose through capitalization.
Incumbent renters are current residents who experience rents, amenities, access, and local development costs.
Potential entrants are not represented in local land-use politics, but may gain when additional housing lets them live in New York rather than take an outside option.

\subsection{Household location choice and entry}

Households choose where to live and work.
A type-\(g\) household choosing residence \(n\) and workplace \(i\) consumes a non-housing good \(c\) and residential floor space \(\ell\).
Direct utility is
\[
    \mathcal{U}_{gni\omega}
    =
    \varepsilon_{gni\omega}
    \frac{B_{gn}(q,L)}{\exp(\kappa \tau_{ni})}
    \left(\frac{c_{gni}}{1-\alpha_g}\right)^{1-\alpha_g}
    \left(\frac{\ell_{gni}}{\alpha_g}\right)^{\alpha_g},
\]
subject to
\[
    c_{gni}+Q_n\ell_{gni}=w_i.
\]
Here \(B_{gn}(q,L)\) is group-specific residential amenity, \(w_i\) is the workplace wage, \(Q_n\) is the residential rent or floor-space price, \(\tau_{ni}\) is commuting time, and \(\alpha_g\) is the housing expenditure share.
The Cobb--Douglas problem implies
\[
    c_{gni}=(1-\alpha_g)w_i,
    \qquad
    Q_n\ell_{gni}=\alpha_g w_i.
\]
Substituting these demands into direct utility yields indirect utility
\[
    U_{gni\omega}
    =
    \varepsilon_{gni\omega}
    \frac{B_{gn}(q,L)w_i}{Q_n^{\alpha_g}\exp(\kappa\tau_{ni})}.
\]

The amenity term contains both positive spatial spillovers and local urban costs:
\[
    B_{gn}(q,L)
    =
    \bar B_{gn}\Omega_n(q,L)^{\eta_g}
    \exp[-\chi_g C_n(q_n,L_n)].
\]
Here \(\bar B_{gn}\) is a baseline amenity, \(\Omega_n(q,L)\) captures positive agglomeration or access spillovers, and \(C_n(q_n,L_n)\) is the true local development cost.
The term \(C_n\) captures congestion, crowding, infrastructure pressure, construction disruption, or other local costs of development.
It is evaluated as an amenity cost in household welfare; it should not also be subtracted as a separate resource cost.

Assume the idiosyncratic term \(\varepsilon_{gni\omega}\) is Fr\'echet with shape parameter \(\theta_g\).
The probability that a type-\(g\) household chooses residence \(n\) and workplace \(i\) is
\[
    \pi_{gni}
    =
    \frac{
    \left[
    B_{gn}(q,L)w_i Q_n^{-\alpha_g}\exp(-\kappa\tau_{ni})
    \right]^{\theta_g}
    }{
    \sum_{n'}\sum_{i'}
    \left[
    B_{gn'}(q,L)w_{i'}Q_{n'}^{-\alpha_g}\exp(-\kappa\tau_{n'i'})
    \right]^{\theta_g}
    }.
\]
The associated inclusive value is
\[
    \Phi_g
    =
    \left[
    \sum_{n'}\sum_{i'}
    \left[
    B_{gn'}(q,L)w_{i'}Q_{n'}^{-\alpha_g}\exp(-\kappa\tau_{n'i'})
    \right]^{\theta_g}
    \right]^{1/\theta_g}.
\]

For incumbent owners and incumbent renters, group populations are fixed in the baseline.
Potential entrants enter the city when the inclusive value of living in New York exceeds the outside option.
A parsimonious entry equation is
\[
    L_N
    =
    \bar L_N
    \left(\frac{\Phi_N}{\bar U_0}\right)^{\psi_N},
\]
where \(\bar U_0\) is the outside option and \(\psi_N\) is the entry elasticity.
This is the channel through which restrictive land use can harm people who are not currently represented in local politics.

\subsection{Housing markets and homeowner assets}

Residential floor-space demand in neighborhood \(n\) is the sum of household demands across groups and workplaces.
Market clearing requires
\[
    \sum_g \sum_i L_g \pi_{gni}\ell_{gni}
    =
    H_n
    =
    \bar H_n+q_n.
\]
Thus, approved housing affects rents, location choices, entry, and ultimately wages and amenities.

Incumbent homeowners differ from renters because they hold claims to local housing assets.
Let \(h_n^O\) denote baseline owner claims on housing assets in neighborhood \(n\).
I treat these claims as fixed at baseline, so that asset-value incidence is evaluated using pre-policy ownership.
A simple capitalization equation is
\[
    P_n = \frac{Q_n}{\iota},
\]
where \(P_n\) is the local housing asset price and \(\iota\) is a capitalization rate.
Homeowner welfare includes both consumption-equivalent utility and capitalization:
\[
    W_O
    =
    \Phi_O
    +
    \sum_n h_n^O(P_n-P_n^0).
\]
Renters and entrants do not receive this asset-value term.

\subsection{Developers and the approval wedge}

Developers produce new housing at physical cost \(K_n(q_n)\), with marginal cost \(K_n'(q_n)\).
Let \(R_n\) denote the marginal market return to new housing in neighborhood \(n\), for example the capitalized value of the marginal unit.
Absent approval frictions, developers would build until
\[
    R_n=K_n'(q_n).
\]
With land-use approval frictions, the supply condition is
\[
    R_n=K_n'(q_n)+\phi_n^s,
    \qquad s\in\{D,C\},
\]
where \(s\) indexes the approval regime and \(\phi_n^s\) is the political or regulatory approval wedge.
The wedge summarizes delay, uncertainty, bargaining, downscaling, required concessions, litigation risk, and outright rejection.

The approval rule below determines the approved quantity \(q_n^s\).
The wedge \(\phi_n^s\) is the implied regulatory cost that rationalizes that approved quantity in the developer's supply condition.
Thus the approval objective and the approval wedge are not two separate supply mechanisms; they are two representations of the same institutional process.

\subsection{Centralized and decentralized approval}

The approval regime determines which privately feasible housing enters the spatial equilibrium.
The social benchmark would evaluate approved housing using true local costs and all affected groups.
The two institutional regimes differ in which benefits and costs they internalize.

\paragraph{Decentralized approval.}
Under decentralized approval, the local decision maker chooses approved housing in neighborhood \(n\) to maximize a politically weighted incumbent objective:
\[
    q_n^D
    \in
    \arg\max_{q_n}\;\widetilde W_n^{inc}(q_n,q_{-n};T_n),
\]
where
\[
    \widetilde W_n^{inc}
    =
    \omega_O(T_n)W_n^O+
    \omega_R(T_n)W_n^R,
    \qquad
    \omega_O'(T_n)>0.
\]
The political weights need not equal welfare weights.
In particular, homeowner exposure can increase the weight placed on incumbent owners.
The decentralized decision maker observes true local costs \(C_n(q_n,L_n)\), but may underweight incumbent renters, potential entrants, and benefits that spill over to other districts.
This is the within-city analogue of an incumbent-controlled planning problem: local incumbents choose growth by trading off the benefits of additional housing against local urban costs, while not fully internalizing the welfare of excluded newcomers.

\paragraph{Centralized approval.}
Under centralized approval, a citywide authority chooses approved housing across neighborhoods:
\[
    q^C
    \in
    \arg\max_q\;\widetilde W^C(q).
\]
A schematic centralized objective is
\[
    \widetilde W^C(q)
    =
    \left[\sum_n W_n^O(q)+\sum_n W_n^R(q)+W^N(q)\right]_{\widetilde C},
\]
where the subscript \(\widetilde C\) indicates that the authority evaluates local development costs using perceived rather than true costs.
The centralized authority internalizes broader benefits, including effects on entrants, rents outside the approving district, and agglomeration.
Its disadvantage is information.
I represent perceived local costs as
\[
    \widetilde C_n(q_n,L_n)
    =
    \psi C_n(q_n,L_n)+\xi_n,
\]
where \(C_n(q_n,L_n)\) is the true local cost of development.
When \(\psi<1\) or \(\xi_n\neq 0\), the central authority mismeasures neighborhood costs.
This is the channel through which decentralized approval can dominate: local officials may observe costs that a citywide authority misses.

\subsection{Three wedges}

The welfare comparison is governed by three wedges.
Let \(W_n^{inc}=W_n^O+W_n^R\) denote local incumbent welfare evaluated with welfare weights, and let \(W^N\) denote entrant welfare.
Suppressing equilibrium feedbacks, the externality wedge is
\[
    \underbrace{
    \frac{\partial W^N}{\partial q_n}
    +
    \sum_{m\neq n}
    \frac{\partial W_m}{\partial q_n}
    }_{\text{externality wedge}}.
\]
This captures benefits to entrants, other districts, citywide rents, labor markets, and agglomeration that a local decision maker may not internalize.

The representation wedge is
\[
    \underbrace{
    \frac{\partial W_n^{inc}}{\partial q_n}
    -
    \frac{\partial \widetilde W_n^{inc}}{\partial q_n}
    }_{\text{representation wedge}}.
\]
This captures the gap between welfare weights and local political weights, especially when incumbent homeowners are overweighted relative to renters and future residents.

The information wedge is
\[
    \underbrace{
    \frac{\partial C_n(q_n,L_n)}{\partial q_n}
    -
    \frac{\partial \widetilde C_n(q_n,L_n)}{\partial q_n}
    }_{\text{information wedge}}.
\]
This captures the part of true local development cost missed by a centralized authority.

Centralized approval is more attractive when the externality and representation wedges are large.
Decentralized approval is more attractive when the information wedge is large.
The model therefore does not assume that either regime is always efficient.
It asks whether the gains from internalizing broader benefits outweigh the loss from imperfect local information.

\subsection{Empirical mapping}

The reduced-form estimates discipline the homeowner-related component of the decentralized approval wedge.
I parameterize this component as
\[
    \phi_{nt}^D
    =
    \phi_n^0+
    \lambda_t\rho T_n+
    u_{nt},
\]
where \(\phi_n^0\) is baseline neighborhood approval difficulty, \(T_n\) is predetermined homeowner exposure, \(\rho\) measures how homeowner exposure affects approval restrictiveness, and \(\lambda_t\) captures the strength of district-based local control over time.

The event-study coefficients are not literal estimates of \(\lambda_t\rho\), because they are quantity responses.
A linearized supply relationship gives
\[
    Y_{nt}
    \approx
    -\eta_q\phi_{nt}^D+\text{controls},
    \qquad
    \beta_t
    \approx
    -\eta_q\lambda_t\rho.
\]
A growing negative homeowner gradient for large-building production therefore disciplines the sign and magnitude of the homeowner-related approval wedge.
Recovering the level of the wedge requires an estimate or calibration of \(\eta_q\), the elasticity of approved housing with respect to the approval wedge.
The reduced-form estimates identify a relative spatial reallocation away from homeowner-heavy districts, not total citywide welfare directly.

\subsection{Welfare counterfactuals}

The model will be used to compare welfare under three regimes.
The first is decentralized approval, in which local political weights and homeowner exposure shape the approval wedge.
The second is centralized approval, in which citywide benefits and entrants receive greater weight but local costs may be mismeasured.
The third is decentralized approval with compensation or pricing.
In that counterfactual, the local decision maker retains local information, but receives a transfer or price that partially internalizes the externality and representation wedges.

Welfare is evaluated separately for incumbent homeowners, incumbent renters, and entrants.
For renters and entrants, welfare is summarized by the group-specific inclusive value.
For homeowners, welfare also includes housing-asset capitalization.
This distinction is important because a regime can reduce local asset values for incumbent owners while increasing welfare for renters and excluded entrants.
